\section{Introduction}

Growing-depth resonance models face two different conditioning problems.  The first is
raw row imbalance: different cutoffs can carry very different numbers of primitive
atoms.  The second is collision accumulation after those rows are assembled in one
coordinate space.  Unit-row normalization removes the first problem by definition,
but it could in principle amplify the second if collision degree grew with depth.

We show that this does not happen in the short-multiplier dilated clock.  The decisive
invariant is not graph degree but coordinate multiplicity.  Exact one-wrap geometry
ensures that no residue bucket contains three prime rows.  Pointwise Cauchy--Schwarz
then gives a Bessel bound of two for the entire normalized family.  The terminology is
standard in frame theory \citep{christensen2016,casazzakutyniok2013}, although the
argument here is a direct one-line coordinate estimate.

The theorem separates two statements that are easy to conflate.  Normalization makes
the operator uniformly bounded as $L$ grows, but it does not force the assembled energy
below its diagonal.  A literal two-row block has ratio $4/3$ in its symmetric mode and
$2/3$ in its antisymmetric mode.  Thus the sign or phase of source coefficients remains
essential.

Finally, normalization changes each row by its norm.  We treat this as a modeling
transform.  Whether the physical arithmetic source admits the same rescaling is a
separate crosswalk theorem, not a consequence of the Bessel estimate.
